\chapter{Inequalities, Large Numbers and Entropy}

\begin{center}\itshape
In $1000$ tosses of a fair coin, $600$ heads is not a rare event for Chebyshev's bound.\\
Hoeffding's bound puts it below $2.1\times 10^{-9}$.
\end{center}

Volume~1, Chapter~10, proves Markov's inequality, Chebyshev's inequality, and the weak law: if $X_1,\ldots,X_n$ are pairwise independent with the same mean $\mu$ and the same finite variance $\sigma^2$, then
\[
 P\bigl(|\overline X_n-\mu|\ge\varepsilon\bigr)\le\frac{\sigma^2}{n\varepsilon^2}.
\]
The bound tends to $0$, so the average converges to $\mu$ in probability. It says nothing about one infinite sequence. This chapter sharpens the tail for bounded variables, proves both Borel--Cantelli lemmas on the coin space of Chapter~1, and deduces a strong law for bounded independent copies. Bernstein's polynomials then turn the weak law into Weierstrass's theorem. Entropy is the same average, applied to $-\log p(X)$.

\section{Chernoff and Hoeffding}

Chebyshev uses only the variance. If the variables are bounded, the exponential moment is finite and gives a much smaller tail.

\begin{lemma}[Hoeffding's lemma]\label{lem:hoeffding}
Let $X$ satisfy $a\le X\le b$ and $\mathbb E X=0$. Then for every real $\lambda$,
\[
 \mathbb E[e^{\lambda X}]\le\exp\Bigl(\frac{\lambda^2(b-a)^2}{8}\Bigr).
\]
\end{lemma}
\begin{proof}
The exponential is convex, so on $[a,b]$ its graph lies under the chord:
\[
 e^{\lambda x}\le\frac{b-x}{b-a}e^{\lambda a}+\frac{x-a}{b-a}e^{\lambda b}.
\]
Since $a\le0\le b$ (otherwise the mean cannot be $0$), both weights below are nonnegative and sum to $1$. Take expectations:
\[
 \mathbb E[e^{\lambda X}]\le\frac{b}{b-a}e^{\lambda a}-\frac{a}{b-a}e^{\lambda b}.
\]
The right-hand side is $\mathbb E[e^{\lambda Z}]$, where $Z$ equals $a$ with probability $b/(b-a)$ and $b$ with probability $-a/(b-a)$. Set $L(\lambda)=\ln\mathbb E[e^{\lambda Z}]$. Then $L(0)=\ln 1=0$. The derivative of a moment generating function at $0$ is the mean, so $L'(0)=\mathbb E Z=0$: the two-point mean is $(ab-ab)/(b-a)=0$. Differentiating again, $L''(\lambda)$ equals the variance of the exponentially tilted copy of $Z$. That copy still takes only the values $a$ and $b$, so it lies in an interval of length $b-a$. For any random variable $W$ in $[a,b]$,
\[
 \operatorname{Var}W=\mathbb E[(W-\mathbb EW)^2]\le\mathbb E\Bigl[\Bigl(W-\frac{a+b}2\Bigr)^2\Bigr]\le\Bigl(\frac{b-a}2\Bigr)^2,
\]
because the midpoint minimises the second moment among constant centres, and $|W-(a+b)/2|\le(b-a)/2$. Thus $L''(\lambda)\le(b-a)^2/4$. Taylor's formula with Lagrange remainder, Volume~1, Chapter~8, gives $L(\lambda)=L''(\xi)\lambda^2/2$ for some $\xi$ between $0$ and $\lambda$. Therefore $L(\lambda)\le\lambda^2(b-a)^2/8$.
\end{proof}

The factor $1/8$ is what the quadratic maximum of $\theta(1-\theta)$ contributes. It is not claimed to be sharp for every law, only for this argument.

\begin{theorem}[Hoeffding's inequality]\label{thm:hoeffding}
Let $X_1,\ldots,X_n$ be independent, with $a_i\le X_i\le b_i$. Set $S=X_1+\cdots+X_n$. Then for every $t>0$,
\[
 P(S-\mathbb E S\ge t)\le\exp\Bigl(-\frac{2t^2}{\sum_{i=1}^n(b_i-a_i)^2}\Bigr),
\]
and the same bound holds for $P(S-\mathbb E S\le -t)$.
\end{theorem}
\begin{proof}
The moment generating function of a sum of independent variables is the product of the factors, because $e^{\lambda X_i}$ are independent and the expectation of a finite product of independent nonnegative variables factors (expand the finite sum that defines each expectation). Lemma~\ref{lem:hoeffding} applied to $X_i-\mathbb E X_i$ gives
\[
 \mathbb E[e^{\lambda(S-\mathbb ES)}]\le\exp\Bigl(\frac{\lambda^2}{8}\sum_i(b_i-a_i)^2\Bigr).
\]
Markov's inequality from Volume~1, applied to the nonnegative variable $e^{\lambda(S-\mathbb ES)}$, yields
\[
 P(S-\mathbb ES\ge t)=P\bigl(e^{\lambda(S-\mathbb ES)}\ge e^{\lambda t}\bigr)
 \le\exp\Bigl(-\lambda t+\frac{\lambda^2}{8}\sum_i(b_i-a_i)^2\Bigr).
\]
Choose $\lambda=4t/\sum_i(b_i-a_i)^2$. The exponent becomes $-2t^2/\sum_i(b_i-a_i)^2$. The lower tail is the same statement for $-X_i$.
\end{proof}

For a fair coin, each summand lies in $\{0,1\}$, so $b_i-a_i=1$ and $\sum(b_i-a_i)^2=n$. If $S$ is the number of heads,
\[
 P(S\ge n/2+t)\le\exp(-2t^2/n).
\]
Take $n=1000$ and $t=100$. Then $600$ heads or more has probability at most $e^{-20}$. Since $e^{-1}<1/2.7<0.371$, one has $e^{-20}<(0.371)^{20}<2.07\times 10^{-9}$. A direct evaluation gives $e^{-20}=2.061\ldots\times 10^{-9}$. The exact binomial tail $\sum_{k=600}^{1000}\binom{1000}{k}2^{-1000}$ equals $1.364\ldots\times 10^{-10}$, computed by summing in logarithmic scale. Hoeffding is larger than the truth by about a factor of fifteen, and Chebyshev is useless: the variance of $S$ is $250$, so
\[
 P(|S-500|\ge100)\le\frac{250}{10000}=\frac1{40}.
\]
One fortieth is a true sentence and a bad estimate.

\begin{figure}[ht]
\centering
\fig{ch04-chernoff}
\caption{Upper bounds on $P(|S_n/n-1/2|\ge 0.1)$ for a fair coin. The charcoal bars are Chebyshev's bound $1/(4n\cdot 0.01)$. The rust bars are Hoeffding's bound $\exp(-2n\cdot 0.01)$. At $n=1000$ the rust bar is $\exp(-20)$, too short to see on this scale; the charcoal bar is still $0.025$.}
\label{fig:chernoff}
\end{figure}

The same exponential estimate specialises to the binomial relative-entropy bound when the summands are Bernoulli, by optimising $\lambda$ against the exact Bernoulli moment $1-p+pe^{\lambda}$ instead of Hoeffding's lemma. For $S\sim\mathrm{Bin}(n,1/2)$ and $a>1/2$,
\[
 P(S\ge an)\le\exp\bigl(-n D(a\|\tfrac12)\bigr),
\]
where $D(a\|\tfrac12)=a\ln(2a)+(1-a)\ln(2(1-a))$. At $a=0.6$ and $n=1000$, $D(0.6\|\tfrac12)=0.020135\ldots$ and the bound is $1.800\ldots\times 10^{-9}$, slightly tighter than $e^{-20}$. The proof is Markov's inequality on $e^{\lambda S}$ and the choice of $\lambda$ that minimises the exponent; it is Exercise~\ref{ex:chernoffbin}.

\section{Borel--Cantelli}

Let $A_1,A_2,\ldots$ be events in the coin space of Chapter~1, each a finite-bit event, so that $P(A_n)$ is a length. Write $A_n\text{ i.o.}$ for the set of outcomes that lie in infinitely many of the $A_n$.

\begin{theorem}[Borel--Cantelli]\label{thm:bc}
\begin{enumerate}[label=(\alph*)]
\item If $\sum_n P(A_n)<\infty$, then $P^*(A_n\text{ i.o.})=0$.
\item If the $A_n$ are mutually independent and $\sum_n P(A_n)=\infty$, then $P^*( (A_n\text{ i.o.})^c )=0$.
\end{enumerate}
\end{theorem}
\begin{proof}
(a) The event $A_n\text{ i.o.}$ is contained in $\bigcup_{n\ge N}A_n$ for every $N$. Proposition~\ref{prop:subadd} gives
\[
 P^*\Bigl(\bigcup_{n\ge N}A_n\Bigr)\le\sum_{n\ge N}P(A_n).
\]
The right-hand side tends to $0$ as $N\to\infty$.

(b) Mutual independence means that every finite subcollection has probability equal to the product of the probabilities, as in Volume~1, Chapter~10. For $0\le x\le1$, $1-x\le e^{-x}$. Hence
\[
 P\Bigl(\bigcap_{n=N}^{N+m}A_n^c\Bigr)=\prod_{n=N}^{N+m}(1-P(A_n))\le\exp\Bigl(-\sum_{n=N}^{N+m}P(A_n)\Bigr).
\]
If the series of probabilities diverges, the right-hand side tends to $0$ as $m\to\infty$. The set where none of $A_N,A_{N+1},\ldots$ occurs is contained in every one of those finite intersections, so its outer content is $0$. That set is the complement of $A_n\text{ i.o.}$ inside the outcomes from index $N$ on. Since $N$ is arbitrary, almost every sequence lies in infinitely many $A_n$.
\end{proof}

Part (a) does not need independence. Part (b) does. Example~\ref{ex:run} is part (b) for the disjoint runs of three heads: $P(A_{3m})=1/8$ and the blocks use distinct digits, so they are independent, and the sum diverges.

\section{The strong law for bounded variables}

The weak law controls each $n$ separately. The strong law says that a single sequence of averages settles.

\begin{theorem}[strong law, bounded case]\label{thm:sllnbound}
Let $X_1,X_2,\ldots$ be independent, identically distributed, and bounded, say $|X_i|\le M$, and suppose each $X_i$ is a function of its own finite block of digits in the coin space of Chapter~1. Let $\mu=\mathbb E[X_1]$. Then
\[
 \frac{X_1+\cdots+X_n}{n}\to\mu
\]
almost surely.
\end{theorem}
\begin{proof}
Write $S_n=X_1+\cdots+X_n$ and $\sigma^2=\operatorname{Var}X_1\le\mathbb E[X_1^2]\le M^2$. Pairwise independence is enough for $\operatorname{Var}S_n=n\sigma^2$, by Volume~1. Chebyshev's inequality gives
\[
 P\bigl(|S_n/n-\mu|\ge\varepsilon\bigr)\le\frac{\sigma^2}{n\varepsilon^2}.
\]
Along $n_k=k^2$ the probabilities sum: $\sum_k\sigma^2/(k^2\varepsilon^2)<\infty$. Each deviation event depends on finitely many digits, because each $X_i$ does and only $n_k$ of them appear. Theorem~\ref{thm:bc}(a) says that almost surely $|S_{n_k}/n_k-\mu|<\varepsilon$ for all large $k$.

Fill the gaps. If $n_k\le n<n_{k+1}$, then $n_{k+1}-n_k=2k+1$ and
\begin{align*}
\Bigl|\frac{S_n}{n}-\mu\Bigr|
&\le\frac{|S_n-S_{n_k}|}{n}+\Bigl|\frac{S_{n_k}}{n_k}-\mu\Bigr|\frac{n_k}{n}+|\mu|\Bigl|1-\frac{n_k}{n}\Bigr|\\
&\le M\frac{2k+1}{k^2}+\Bigl|\frac{S_{n_k}}{n_k}-\mu\Bigr|+|\mu|\frac{2k+1}{k^2}.
\end{align*}
The deterministic terms tend to $0$. Almost surely the middle term is eventually smaller than $\varepsilon$. So the limsup of $|S_n/n-\mu|$ is at most $\varepsilon$. Let $\varepsilon\downarrow0$ through the sequence $1/m$. A countable union of outer-content-zero sets still has outer content zero, by Proposition~\ref{prop:subadd}.
\end{proof}

Any i.i.d.\ sequence with finitely many values can be realised on disjoint blocks of fair bits, so the theorem covers coin tosses, dice, and every finite alphabet. A bounded variable with a density is built in Chapter~5 from a uniform digit block by the inverse distribution function. The algebra above does not change; only the construction of the copies is postponed.

The hypothesis $|X_i|\le M$ is used to control the gaps between $k^2$ and $(k+1)^2$. A strong law with only $\mathbb E|X_1|<\infty$, and no bound, is proved in Volume~11. It is not a step in this volume.

\section{Bernstein's proof of Weierstrass's theorem}

Let $f$ be continuous on $[0,1]$. The Bernstein polynomial of degree $n$ is
\[
 B_n(f)(x)=\sum_{k=0}^n f\Bigl(\frac kn\Bigr)\binom{n}{k}x^k(1-x)^{n-k}.
\]
If $S$ is binomial with parameters $n$ and $x$, then $B_n(f)(x)=\mathbb E[f(S/n)]$.

\begin{theorem}[Weierstrass, via Bernstein]\label{thm:weierstrass}
If $f$ is continuous on $[0,1]$, then $B_n(f)\to f$ uniformly on $[0,1]$.
\end{theorem}
\begin{proof}
Continuity on a closed bounded interval gives uniform continuity and a bound $|f|\le M$, both from Volume~1, Chapter~3. Fix $\varepsilon>0$ and choose $\delta>0$ so that $|y-x|<\delta$ implies $|f(y)-f(x)|<\varepsilon$. Then
\begin{align*}
|B_n(f)(x)-f(x)|
&=\bigl|\mathbb E[f(S/n)-f(x)]\bigr|\\
&\le\varepsilon\,P(|S/n-x|<\delta)+2M\,P(|S/n-x|\ge\delta).
\end{align*}
The binomial variance is $nx(1-x)\le n/4$, so Chebyshev's inequality yields $P(|S/n-x|\ge\delta)\le 1/(4n\delta^2)$. For every $n>M/(2\varepsilon\delta^2)$ the second term is smaller than $\varepsilon$, uniformly in $x$. The first term is at most $\varepsilon$.
\end{proof}

Polynomials are therefore dense in the continuous functions on a segment, in the uniform norm. Volume~6 cites this theorem when it approximates a continuous weight. The proof above is the one that citation points to.

\section{Entropy of a finite alphabet}

Let $X$ take finitely many values, with $P(X=a)=p_a>0$. The \emph{Shannon entropy} in bits is
\[
 H(X)=\sum_a p_a\log_2\frac1{p_a}.
\]
The sum is finite. Changing the base of the logarithm multiplies $H$ by a constant; bits are the choice $\log_2$. A fair coin has $H=1$. A coin with $P(\mathrm{heads})=0.1$ has
\[
 H=0.1\log_2 10+0.9\log_2(10/9)=0.46899\ldots.
\]

\begin{proposition}[Gibbs]\label{prop:gibbs}
If $X$ takes at most $m$ values, then $H(X)\le\log_2 m$, with equality if and only if $X$ is uniform on $m$ values.
\end{proposition}
\begin{proof}
The inequality $\ln u\le u-1$ for $u>0$, from Volume~1, Chapter~9, gives, for any probability masses $q_a>0$ on the same values,
\[
 \sum_a p_a\ln\frac{q_a}{p_a}\le\sum_a p_a\Bigl(\frac{q_a}{p_a}-1\Bigr)=\sum_a q_a-\sum_a p_a=0.
\]
Thus $\sum p_a\log_2(1/p_a)\le\sum p_a\log_2(1/q_a)$. Take $q_a=1/m$. Equality holds when $q_a=p_a$ for every $a$, because $\ln u=u-1$ only at $u=1$.
\end{proof}

For i.i.d.\ copies $X_1,X_2,\ldots$ with the same finite law, the probability of a particular string is the product of the symbol probabilities, so
\[
 -\frac1n\log_2 P(X_1=a_1,\ldots,X_n=a_n)=\frac1n\sum_{i=1}^n\log_2\frac1{p_{a_i}}.
\]
The summands are i.i.d.\ and bounded, because every $p_a$ in the support is positive and the alphabet is finite. Theorem~\ref{thm:sllnbound} says the average converges almost surely to $\mathbb E[\log_2(1/p_X)]=H(X)$.

\begin{theorem}[typical sequences]\label{thm:aep}
Let $\varepsilon>0$. The set of strings of length $n$ whose probability $p$ satisfies
\[
 2^{-n(H(X)+\varepsilon)}\le p\le 2^{-n(H(X)-\varepsilon)}
\]
has probability tending to $1$. Every such string has probability at most $2^{-n(H(X)-\varepsilon)}$, so the set contains at most $2^{n(H(X)+\varepsilon)}$ strings once $n$ is large enough that its probability exceeds $1/2$.
\end{theorem}
\begin{proof}
The almost-sure convergence just recorded implies convergence in probability of the same averages, which is the first sentence. If a set of probability greater than $1/2$ consists of points each of mass at most $2^{-n(H-\varepsilon)}$, then the number of points is at least one, and at most the reciprocal of that mass times $1$, hence at most $2^{n(H+\varepsilon)}$ after adjusting $\varepsilon$ by the factor coming from the lower bound in the display. More directly: each listed string has mass at least $2^{-n(H+\varepsilon)}$, and the masses sum to at most $1$, so there are at most $2^{n(H+\varepsilon)}$ of them.
\end{proof}

The count $2^{nH}$ is the size of the set of strings one actually sees. The whole space may be vastly larger. That is the content of the entropy.

The same $H$, with natural logarithms, was used by Gilmer in 2022 to bound the largest frequency in a finite union-closed family. The comparison that produces the constant $(1-\ln 2)/2$ is not proved in this volume. It is not used later in the series by label, and it is left out rather than sketched.

Chapter~5 leaves countable values. A uniform point of $[0,1]$ has no mass at any single number, and an expectation becomes an integral.

\section{Exercises}

\begin{exercise}[$\star$]\label{ex:markovsharp}
Let $a>0$ and $0<p\le1$. Construct a nonnegative random variable with $\mathbb E X=a$ and $P(X\ge a)=p$, or show that the second demand cannot be met unless $p\le1$. Then show that Markov's inequality is attained by a variable that takes only the values $0$ and $a$.
\end{exercise}

\begin{exercise}[$\star\star$]\label{ex:chernoffbin}
Let $S$ be binomial$(n,1/2)$ and let $a\in(1/2,1)$. Prove
\[
 P(S\ge an)\le\exp\bigl(-n D\bigr),\qquad D=a\ln(2a)+(1-a)\ln\bigl(2(1-a)\bigr),
\]
by applying Markov's inequality to $e^{\lambda S}$ and minimising the resulting exponent in $\lambda>0$.
\end{exercise}

\begin{exercise}[$\star$]\label{ex:bcfinite}
Let $A_n$ be independent events with $P(A_n)=1/n$. Show that almost every outcome lies in infinitely many $A_n$. If instead $P(A_n)=1/n^2$, show that almost every outcome lies in only finitely many.
\end{exercise}

\begin{exercise}[$\star\star$]\label{ex:bernstein2}
Let $f(x)=x^2$ on $[0,1]$. Compute $B_2(f)$ and $B_n(f)$ in closed form, and check that $B_n(f)(x)-x^2=x(1-x)/n$.
\end{exercise}

\begin{exercise}[$\star\star$]\label{ex:entropymax}
For a Bernoulli$(p)$ variable, $0<p<1$, show that $H(p)=-p\log_2 p-(1-p)\log_2(1-p)$ has a maximum at $p=1/2$, and that the maximum value is $1$.
\end{exercise}

\begin{exercise}[$\star\star$]\label{ex:gapfill}
In the proof of Theorem~\ref{thm:sllnbound}, the subsequence was $n_k=k^2$. Show that $n_k=k$ does not make $\sum_k P(|S_{n_k}/n_k-\mu|\ge\varepsilon)$ converge if one only has Chebyshev's bound, and that $n_k=\lceil k^{1+\delta}\rceil$ does, for every $\delta>0$.
\end{exercise}

\begin{exercise}[$\star\star\star$]\label{ex:runs}
Let $R_n$ be the longest run of heads in $n$ fair tosses. Show that $R_n/\log_2 n\to 1$ in probability. A complete almost-sure statement is not required.
\end{exercise}

\section{Hints}

\begin{hint}{ex:markovsharp}
Put mass $p$ at $a/p$ only if you want expectation $a$ and a threshold $a$. For equality in Markov, two-point variables suffice.
\end{hint}
\begin{hint}{ex:chernoffbin}
$\mathbb E[e^{\lambda S}]=( (e^{\lambda}+1)/2 )^n$. Take logarithms and set the derivative in $\lambda$ to zero.
\end{hint}
\begin{hint}{ex:bcfinite}
Compare $\sum 1/n$ and $\sum 1/n^2$ with the integral test from Volume~1, and apply Theorem~\ref{thm:bc}. The events have to be realised as independent finite-bit events; disjoint blocks of bits can do that for these probabilities only approximately. Use the theorem as stated for any independent family already constructed, and note that $P=1/n$ may be replaced by $c/n$.
\end{hint}
\begin{hint}{ex:bernstein2}
$B_n(f)(x)=\mathbb E[(S/n)^2]=\operatorname{Var}(S/n)+(\mathbb E[S/n])^2$.
\end{hint}
\begin{hint}{ex:entropymax}
Differentiate. The second derivative of $-p\ln p-(1-p)\ln(1-p)$ is negative.
\end{hint}
\begin{hint}{ex:gapfill}
Chebyshev contributes a factor $1/n_k$. The series $\sum k^{-1-\delta}$ converges.
\end{hint}
\begin{hint}{ex:runs}
If $r=\lfloor(1+\varepsilon)\log_2 n\rfloor$, there are at most $n$ windows of length $r$, each with probability $2^{-r}$. If $r=\lfloor(1-\varepsilon)\log_2 n\rfloor$, look at about $n/r$ disjoint windows.
\end{hint}

\section{Solutions}

\begin{solution}{ex:markovsharp}
\idea Markov compares $X$ with $a$ times an indicator. Equality holds when $X$ equals that multiple.
\writeup If $P(X\ge a)=p$ and $X\ge0$ with $\mathbb E X=a$, Markov says $p\le1$, so $p>1$ is impossible. For equality at a prescribed $p\in(0,1]$, let $X=0$ with probability $1-p$ and $X=a/p$ with probability $p$. Then $\mathbb E X=a$ and $P(X\ge a)=p$, but $P(X\ge a)=\mathbb E X/a$ only when $p=1$. The equality case of the proof is the variable that is $0$ or $a$: then $a\mathbf 1_{X\ge a}=X$ at every outcome, so taking expectations gives equality.
\end{solution}

\begin{solution}{ex:chernoffbin}
\idea The Bernoulli moment generating function is elementary, and the best $\lambda$ is the one that matches the threshold $a$ to the tilted mean.
\writeup $\mathbb E[e^{\lambda S}]=(1/2+e^{\lambda}/2)^n$. Markov's inequality gives
\[
 P(S\ge an)\le e^{-\lambda an}\Bigl(\frac{1+e^{\lambda}}{2}\Bigr)^n
 =\exp\Bigl(-n\bigl(\lambda a-\ln(1+e^{\lambda})+\ln 2\bigr)\Bigr).
\]
Let $\phi(\lambda)=\lambda a-\ln(1+e^{\lambda})$. Then $\phi'(\lambda)=a-e^{\lambda}/(1+e^{\lambda})$. Set the derivative to zero: $e^{\lambda}/(1+e^{\lambda})=a$, so $e^{\lambda}=a/(1-a)$ and $\lambda=\ln(a/(1-a))>0$ because $a>1/2$. Substitute:
\[
 \phi(\lambda)=a\ln\frac{a}{1-a}-\ln\Bigl(\frac{1}{1-a}\Bigr)=a\ln a+(1-a)\ln(1-a).
\]
The exponent is $-n(\phi(\lambda)+\ln 2)=-n\bigl(a\ln(2a)+(1-a)\ln(2(1-a))\bigr)$.
\end{solution}

\begin{solution}{ex:bcfinite}
\idea The harmonic series diverges and the square series converges. Independence is required only for the divergent half.
\writeup The integral test of Volume~1 gives $\sum_{n\ge1}1/n=\infty$ and $\sum 1/n^2<\infty$. Theorem~\ref{thm:bc}(a) gives the convergent case, with no independence. For $P(A_n)=1/n$, part (b) needs an independent realisation. Disjoint blocks of $n$ fair bits can produce an event of probability $2^{-n}$, not $1/n$. Realise instead a sequence of independent events of probability $1/2$: the single bits $b_n$. Those satisfy $\sum P=\infty$ and fall under part (b), which is the intent of the exercise for the divergent regime. For probabilities exactly $1/n$, take independent geometric waiting times on the countable space of Volume~1, one for each $n$, on a countable product of countable sets, which is countable, and let $A_n$ be the event that the $n$th wait exceeds a threshold chosen so that the probability is $1/n$. Finite subcollections are independent by construction of the product masses. Part (b) applies and the sum diverges, so almost every outcome lies in infinitely many $A_n$.
\end{solution}

\begin{solution}{ex:bernstein2}
\idea The second moment of a binomial random variable is $np(1-p)+(np)^2$.
\writeup With $x$ as the success probability, $\mathbb E[S]=nx$ and $\mathbb E[S^2]=nx(1-x)+n^2x^2$. Hence
\[
 B_n(f)(x)=\mathbb E[(S/n)^2]=x^2+\frac{x(1-x)}{n}.
\]
For $n=2$, $B_2(f)(x)=x^2+x(1-x)/2$. The difference is $x(1-x)/n$, which tends uniformly to $0$ because $x(1-x)\le1/4$.
\end{solution}

\begin{solution}{ex:entropymax}
\idea The binary entropy is concave, and its derivative vanishes only at $1/2$.
\writeup Work in natural units and convert at the end. Let $h(p)=-p\ln p-(1-p)\ln(1-p)$. Then
\[
 h'(p)=-\ln p+\ln(1-p)=\ln\frac{1-p}{p},\qquad h''(p)=-\frac1p-\frac1{1-p}<0.
\]
So $h'$ vanishes only at $p=1/2$, and $h$ has a maximum there. Now $h(1/2)=\ln 2$, and dividing by $\ln 2$ converts to bits, giving $H(1/2)=1$. For every other $p\in(0,1)$, $H(p)<1$.
\end{solution}

\begin{solution}{ex:gapfill}
\idea Chebyshev contributes one factor of $1/n_k$. The exponent of $k$ has to be strictly larger than $1$.
\writeup If $n_k=k$, the Chebyshev bound is $\sigma^2/(k\varepsilon^2)$, and $\sum 1/k$ diverges, so this particular estimate does not feed Theorem~\ref{thm:bc}(a). The bound may still go to zero; the series of bounds does not converge. If $n_k=\lceil k^{1+\delta}\rceil$, then $1/n_k\le k^{-1-\delta}$, and $\sum k^{-1-\delta}$ converges by comparison with the integral of $x^{-1-\delta}$.
\end{solution}

\begin{solution}{ex:runs}
\idea A window of length $(1+\varepsilon)\log_2 n$ is too rare for $n$ of them to succeed. A window of length $(1-\varepsilon)\log_2 n$ is common enough that disjoint copies succeed.
\writeup Let $r=\lfloor(1+\varepsilon)\log_2 n\rfloor$. There are at most $n$ contiguous blocks of $r$ tosses. Each has probability $2^{-r}\le 2\cdot n^{-(1+\varepsilon)}$ once $n$ is large, because $2^{-r}\le 2^{-(1+\varepsilon)\log_2 n+1}=2 n^{-(1+\varepsilon)}$. The union bound gives probability at most $2n\cdot n^{-(1+\varepsilon)}=2n^{-\varepsilon}\to0$. So $P(R_n\ge(1+\varepsilon)\log_2 n)\to0$.

Let $s=\lfloor(1-\varepsilon)\log_2 n\rfloor$ and look at $m=\lfloor n/(s+1)\rfloor$ disjoint blocks. The probability that a given block is all heads is $2^{-s}\ge n^{-(1-\varepsilon)}/2$ for large $n$. The probability that none of them is all heads is at most $(1-c n^{-(1-\varepsilon)})^m$. Here $m$ is about $n/((1-\varepsilon)\log_2 n)$, so $m n^{-(1-\varepsilon)}$ tends to infinity. Thus the product tends to $0$, and $P(R_n<s)\to0$. Combining the two estimates gives $R_n/\log_2 n\to 1$ in probability.
\end{solution}

\section*{Sources and further reading}

\begin{enumerate}
\item J.~Bernoulli, \emph{Ars Conjectandi}, Basel, 1713. The weak law for frequencies, proved there by direct binomial estimates.
\item P.~L.~Chebyshev, Des valeurs moyennes, \emph{J.\ Math.\ Pures Appl.} (2) \textbf{12} (1867), 177--184. The inequality Volume~1 uses.
\item W.~Hoeffding, Probability inequalities for sums of bounded random variables, \emph{J.\ Amer.\ Statist.\ Assoc.} \textbf{58} (1963), 13--30. Theorem~\ref{thm:hoeffding}.
\item \'E.~Borel, Les probabilit\'es d\'enombrables et leurs applications arithm\'etiques, \emph{Rend.\ Circ.\ Mat.\ Palermo} \textbf{27} (1909), 247--271. The first lemma, for independent digits. F.~P.~Cantelli, Sulla probabilit\`a come limite della frequenza, \emph{Atti Accad.\ Naz.\ Lincei} \textbf{26} (1917), 39--45, removed the independence hypothesis from that half.
\item S.~Bernstein, D\'emonstration du th\'eor\`eme de Weierstrass fond\'ee sur le calcul des probabilit\'es, \emph{Comm.\ Soc.\ Math.\ Kharkov} \textbf{13} (1912), 1--2. Theorem~\ref{thm:weierstrass}.
\item C.~E.~Shannon, A mathematical theory of communication, \emph{Bell System Tech.\ J.} \textbf{27} (1948), 379--423 and 623--656. The entropy of a finite alphabet.
\end{enumerate}
